Protein binding pocket identification is crucial for drug discovery, but conventional methods struggle with accuracy and physical plausibility. We present PhyGNN, a physics-informed graph neural network that integrates Hamiltonian mechanics constraints to improve pocket detection. By representing proteins as graphs with physics-enhanced features and incorporating energy conservation, bond length, and electrostatic constraints, PhyGNN achieves state-of-the-art performance with an F1 score of 0.547 on the PDBbind dataset. This represents a 5.2\% improvement over established geometric methods (FPOCKET F1=0.520) and a 77.7\% improvement over non-physics GNNs. Through ablation studies, we demonstrate that electrostatic interactions contribute most significantly to performance improvements. Application to therapeutic targets including STAT3 (cancer immunotherapy), KRAS G12C (oncology), and LRRK2 (Parkinson's disease) reveals previously unidentified potential binding sites. PhyGNN provides a framework for physics-aware machine learning in structural biology, with implications for cryptic pocket discovery and drug design against currently "undruggable" targets.
